% !TEX root = ../main.tex

\chapter{吹熄测量方法的建立}

\section{引言}

在热流工况下，如何建立可重复、可量化的横向射流火焰吹熄判定方法。准确判定火焰吹熄是开展后续收缩流道影响研究的前提。为此，本章从同温工况筛选、二级燃料优化、信号采集改进和占空比判据建立四个方面，系统阐述吹熄测量方法的建立过程。

\section{同温工况筛选}

\subsection{筛选的必要性}

在开展二级射流吹熄实验之前，必须先把一级燃烧的条件确定下来。一级燃烧室的作用是提供稳定的高温主流，但改变空气流量时，主流温度和速度会同时变化——如果不先确定“哪些工况温度相同”，后续就无法判断二级火焰的变化是由温度引起的还是由速度引起的。因此，同温工况筛选是后续所有对比实验的前提。

\subsection{筛选方法与流程}

筛选流程分为三个步骤：

第一步，初筛。固定一级甲烷流量，逐步改变一级空气流量（300-900 SLM范围），记录二级燃烧室入口B型热电偶的温度。每个工况点稳定至少30秒后记录读数。

第二步，分组。将数据按温度分组，容差设定为 \(\pm 15^{\circ}C\) ，将温度相近的工况归为同一组。

第三步，验证。将初筛的工况各稳定燃烧3分钟，排除温度波动超限的点。

\begin{figure}[H]
	\centering
	\includegraphics[width=0.8\textwidth]{fig5}
	\caption{同温工况筛选结果}
	\label{fig:isothermal_conditions}
\end{figure}

\subsection{筛选结果}

通过多轮筛选，确定了一组同温工况：

\begin{table}[!htbp]
	\centering
	\caption{同温工况组合}
	\begin{tabular}{c|c|c}
		\hline
		一级空气流量(SLM) & 甲烷流量(SLM) & 温度(K) \\
		\hline
		360 & 32 & 1342 \\
		455 & 31 & 1344 \\
		510 & 30 & 1343 \\
		560 & 31 & 1342 \\
		600 & 30 & 1341 \\
		650 & 32 & 1343 \\
		720 & 32 & 1341 \\
		770 & 33 & 1342 \\
		825 & 34 & 1342 \\
		870 & 35 & 1343 \\
		930 & 36 & 1342 \\
		970 & 37 & 1342 \\
		\hline
	\end{tabular}
\end{table}

以上工况主流温度均在1342K左右（容差 \(\pm 5\mathrm{K}\) ），主流速度随空气流量增加而单调递增。该工况表覆盖了从稳定燃烧到接近吹熄的完整速度范围，为后续二级射流对比实验提供了可对比的基础条件。

\section{二级射流燃料优化}

\subsection{甲烷燃料的局限性}

初期使用甲烷作为二级燃料，蓝色火焰在高温主流背景下对比度极低。低射流流量时肉眼几乎无法辨认火焰是否存在，吹熄临界判定主观性强。尝试调节相机参数，改善有限。

\subsection{乙烯燃料的优化效果}

将二级燃料从甲烷更换为乙烯后，综合发光信号有所增强，在部分工况下可观测性相比甲烷有一定改善。更换燃料涉及气路改造、流量计重新标定、安全操作流程更新等环节。

\begin{figure}[H]
	\centering
	\includegraphics[width=0.8\textwidth]{fig6}
	\caption{乙烯二级射流火焰}
	\label{fig:flame_visibility}
\end{figure}

但在接近吹熄的工况下，乙烯火焰同样存在可见度不足的问题——这也正是后续需要引入紫外探测器进行量化测量的根本原因。

\section{火焰信号采集改进}

使用YJH-1E紫外火焰探测器（4-20mA输出）搭配模拟采集模块监测火焰信号。初步使用配套软件采集时，发现数据刷新速度慢，无法记录火焰的快速变化。当火焰进入吹熄前的“频闪”状态时，探测器输出基本不变。

为排查问题根源，使用示波器直接测量探测器的4-20mA输出端，发现探测器本身对火焰变化的响应速度远快于配套软件的采集存储频率——瓶颈在于采集模块的软件端。

基于这一判断，放弃原配套采集方案，改用Python编写串口读取程序，直接通过RS-485接口读取探测器数据。优化后，有效数据存储频率从原系统的5Hz提升至约20Hz，实现了4倍的采集速度提升。20Hz采样虽不足以进行波形细节分析，但已足以反映吹熄前信号的起伏趋势。

\begin{figure}[H]
	\centering
	\includegraphics[width=0.8\textwidth]{fig7}
	\caption{20Hz采样时稳定存在火焰信号}
	\label{fig:signal_fluctuation}
\end{figure}

\section{占空比判据的建立}

\subsection{占空比的定义}

将探测器电流信号二值化——电流大于4.2mA判为“有火”，电流小于等于4.2mA判为“无火”。4.2mA的阈值对应探测器的暗电流输出（4mA）加0.2mA的抗噪声裕量。占空比即单位时间内“有火”状态的时间占总时间的比例，反映了火焰在统计时间内的“存在概率”。

\subsection{四级吹熄判据}

结合实验现象，建立了四级判据：

第一阶段：持续燃烧阶段。火焰整体维持在喷嘴出口附近，占空比接近1.00。火焰面在湍流作用下会发生抖动和局部变形，偶尔出现瞬时信号低于阈值，但火焰整体连续性不受影响。

第二阶段：近场频闪。火焰总体位置仍在喷嘴出口附近，占空比0.5-0.8附近，已无法保持持续燃烧状态，开始呈现频繁的亮暗交替。火焰处于“勉强维持”的边缘状态。

第三阶段：下游偶发复燃。火焰根部已无法维持在喷嘴出口，被吹至下游较远位置，占空比0.1-0.3附近。绝大部分时间内无火焰信号，偶有下游短暂复燃但很快再次熄灭。

第四阶段：完全吹熄。火焰信号完全消失，占空比趋近于0，下游也不再出现复燃。

\begin{table}[!htbp]
	\centering
	\caption{四级吹熄判据体系}
	\begin{tabular}{c|c|c}
		\hline
		阶段 & 占空比范围 & 现象 \\
		\hline
		Ⅰ级持续燃烧 & \(\approx 1.00\) & 火焰连续稳定，偶有抖动 \\
		Ⅱ级近场频闪 & \(0.5-0.8\) & 频繁亮暗交替，勉强维持 \\
		Ⅲ级下游偶发复燃 & \(0.1-0.3\) & 长时间无火，偶发短暂燃烧 \\
		Ⅳ级完全吹熄 & \(\approx 0\) & 火焰信号完全消失 \\
		\hline
	\end{tabular}
\end{table}

上述阶段划分及对应的占空比范围是基于当前有限工况下的实验观测和数据分析得出的初步结果，后续可结合更完整的实验数据进行进一步验证和完善。

\section{本章小结}

本章完成了以下工作：

（1）通过多轮筛选确定多个同温工况（温度约 \(1342\pm 5\mathrm{K}\) ），为后续吹熄实验打好基础。

（2）将二级燃料从甲烷优化为乙烯，优化了火焰可见度不足的问题，为后续分析记录火焰形态做准备。

（3）通过Python高频读取将探测器采样频率从5Hz提升至20Hz，成功通过火焰探测器观察到低频的吹熄前频闪。

（4）建立了基于占空比的四级吹熄判据体系，使火焰稳定性从主观判断转变为可重复计算的量化指标。